\documentclass[areasetadvanced]{scrartcl}

\usepackage[utf8]{inputenc}
\usepackage[T2A]{fontenc}
\usepackage[english,russian]{babel}

\usepackage[footskip=1cm,left=25mm, right=15mm, top=20mm, bottom=20mm]{geometry}
\usepackage{setspace}
\usepackage{amsmath, amssymb} 
\usepackage{graphicx}
\usepackage{tikz}
\usetikzlibrary{arrows.meta}
\usepackage{float}
\usepackage{dashrule}
\usepackage{fancyhdr} 
\usepackage{hyperref} 
\usepackage{parskip}
\usepackage{textcomp, enumitem}
\usepackage{indentfirst}
\usepackage{graphicx}
\usepackage{algorithm}
\usepackage{algpseudocode}
\usepackage{array} 
\usepackage{geometry}
\usepackage{afterpage}
\usepackage{minted}
\setcounter{secnumdepth}{3} 
\setcounter{tocdepth}{3}    
\usepackage{listings} 

\newcommand{\icon}[1]{\includegraphics[height=1.2em]{\#1}}

\tikzstyle{block} = [rectangle, rounded corners, minimum width=3cm, minimum height=1cm, text centered, draw=black, fill=lightgray]

\setkomafont{sectioning}{\normalfont\bfseries} 
\setkomafont{section}{\normalfont\Large\bfseries}
\setkomafont{subsection}{\normalfont\large\bfseries}
\setkomafont{subsubsection}{\normalfont\large\bfseries}
\setkomafont{paragraph}{\normalfont\large\bfseries} 

\lstset{
  language=Haskell,
  basicstyle=\ttfamily\small,
  keywordstyle=\color{blue}\bfseries,
  stringstyle=\color{red},
  commentstyle=\color{green!70!black},
  numbers=left,
  numberstyle=\tiny,
  stepnumber=1,
  numbersep=10pt,
  showstringspaces=false,
  breaklines=true,
  frame=single
}

\lstdefinelanguage{Lua}{
    keywords={function, end, if, then, else, elseif, for, while, do, repeat, until, break, return, local, and, or, not, true, false, nil},
    keywordstyle=\color{blue}\bfseries,
    stringstyle=\color{red},
    commentstyle=\color{green!70!black},
    morestring=[s]{"}{"},
    morestring=[s]{'}{'},
    morecomment=[l]{--},
    morecomment=[s]{--[[}{]]},
    basicstyle=\ttfamily\small,
    numbers=left,
    numberstyle=\tiny,
    stepnumber=1,
    numbersep=10pt,
    showstringspaces=false,
    breaklines=true,
    frame=single
}

\setlength{\parindent}{1.25cm}
\setcounter{tocdepth}{3}
\begin{document}
\sloppy
	\thispagestyle{empty}
	\begin{center}
		\large{МИНОБРНАУКИ РОССИИ} \par
		\vspace{0.3cm}
		\normalsize
		{ФЕДЕРАЛЬНОЕ ГОСУДАРСТВЕННОЕ АВТОНОМНОЕ ОБРАЗОВАТЕЛЬНОЕ УЧРЕЖДЕНИЕ ВЫСШЕГО ОБРАЗОВАНИЯ} \par
		\vspace{0.3cm}
		\textbf{\guillemotleft САНКТ-ПЕТЕРБУРГСКИЙ ПОЛИТЕХНИЧЕСКИЙ}
		\textbf{УНИВЕРСИТЕТ ПЕТРА ВЕЛИКОГО\guillemotright} \par
		\vspace{0.3cm}
		{Институт компьютерных наук и кибербезопасности}\par
		{Высшая школа технологий искусственного интеллекта}\par
	\end{center}
	\vfill
	\begin{center}
        \par
		{\huge Методы исследовательской работы}\par
        \Large {Использование настраиваемых персон и аватаров в LLM-чат-ботах:
как кастомизация образа ИИ влияет на самораскрытие, регуляцию аффекта
и переживание «цифровой близости».}\par
	\end{center}
	\vfill
	\begin{flushleft}
		Студент: \hspace{1.8cm} \rule[0pt]{2.5cm}{0.5pt}\hfill Салимли Айзек Мухтар Оглы\par
		\vspace{1.5cm}
		Преподаватель: \hspace{0.55cm} \rule[0pt]{2.5cm}{0.5pt}\hfill  Заборовский Владимир Сергеевич
	\end{flushleft}
	\vspace{0.5cm}
	\begin{flushright}
		\guillemotleft \rule[0pt]{0.8cm}{0.5pt}\guillemotright \rule[0pt]{2cm}{0.5pt} 20\rule[0pt]{0.5cm}{0.5pt} г.
	\end{flushright}
	\vfill
	\begin{center}
		Санкт-Петербург, 2026
	\end{center}
	\newpage
	\tableofcontents
\newpage
\section*{Введение}
\addcontentsline{toc}{section}{Введение}

\subsection*{Актуальность темы}
Рост популярности ИИ-ассистентов, стремление к персонализации, «одиночество в цифровую эпоху».

\subsection*{Проблема исследования}
Как техническая возможность настройки внешности и характера ИИ меняет психологический опыт взаимодействия.

\subsection*{Объект и предмет исследования}
% Текст раздела.

\subsection*{Цель и гипотезы работы}
% Текст раздела.

\subsection*{Научная новизна и практическая значимость}
% Текст раздела.

\newpage
\section{Теоретические основы и контекст исследования}

\subsection{Эволюция человеко-машинного взаимодействия: от командной строки к антропоморфным агентам}
% Текст раздела.

\subsection{Психология общения с неодушевленным объектом}
Эффект Элизы (склонность приписывать машинам человеческие черты). Парасоциальные отношения с медиаперсонажами и ботами.
% Текст раздела.

\subsection{Ключевые понятия исследования}
\begin{itemize}
\item \textbf{Самораскрытие (Self-disclosure):} теории, риски и преимущества.
\item \textbf{Регуляция аффекта (Affect regulation):} как общение помогает справляться со стрессом.
\item \textbf{Цифровая близость (Digital intimacy):} определение, компоненты, отличие от реальной близости.
\end{itemize}

\newpage
\section{Обзор и анализ современных инструментов: кастомизация персон и аватаров}

\subsection{Концепция «персонажа» в современных LLM: от безличного помощника к собеседнику}
% Текст раздела.

\subsection{Экосистема OpenAI (ChatGPT)}
Возможности кастомизации через инструкции (Custom instructions). Генерация аватаров (DALL-E интеграция) и создание образа. Специализированные GPTs: создание персон с уникальным промптом и знаниями.
% Текст раздела.

\subsection{Платформа DeepSeek}
Доступный функционал кастомизации (стиль ответов, память). Специфика взаимодействия и возможности настройки «характера» в свободном режиме.
% Текст раздела.

\subsection{Семейство Grok AI (xAI)}
Встроенная персона: «бунтарь», «юморист». Анализ заданной, а не настраиваемой персоны.
% Текст раздела.

\subsubsection{Grok Companion (виртуальный партнер)}
Фокус на эмоциональную связь и более глубокое взаимодействие. Анализ того, как разработчики проектируют «близость» и какие инструменты кастомизации предлагаются пользователю для усиления этого эффекта.
% Текст раздела.

\subsection{Типология сценариев взаимодействия с ИИ-компаньонами}
Анализ типичных контекстов использования и поведение систем в различных ситуациях.

\subsubsection{Сценарий установления контакта и приветствия}
% Текст раздела.

\subsubsection{Партнёрские и романтические сценарии}
% Текст раздела.

\subsubsection{Реакция на агрессивное поведение пользователя: грубость и оскорбления}
% Текст раздела.

\subsubsection{Ситуации ревности}
% Текст раздела.

\subsubsection{Конфликтные и эскалационные ситуации}
% Текст раздела.

\subsection{Психологическая поддержка и этические границы}
Возможности и ограничения ИИ-компаньонов в контексте психического здоровья пользователей.

\subsubsection{Психологическая помощь и эмпатическая поддержка}
% Текст раздела.

\subsubsection{Рекомендации и советы: возможности и ограничения}
% Текст раздела.

\subsubsection{Разграничение компетенций: ситуации, требующие помощи психиатра}
Анализ того, даёт ли система совет о необходимости обращения к специалисту (психиатру) и насколько корректно определяет границы применимости.
% Текст раздела.

\subsubsection{Обработка интимных откровений и конфиденциальность}
% Текст раздела.

\subsection{Сравнительный анализ}
Таблица сравнения платформ по критериям (глубина кастомизации, тип аватара, ориентация на эмоциональный отклик).
% Текст раздела.

\newpage
\section{Методология и дизайн эмпирического исследования}

\subsection{Постановка задачи и этапы исследования}
% Текст раздела.

\subsection{Описание выборки: респонденты, критерии отбора}
% Текст раздела.

\subsection{Методы сбора данных}
\begin{itemize}
\item Шкалы самораскрытия (например, адаптация опросника Jourard).
\item Шкалы оценки аффективного состояния (PANAS или аналоги).
\item Полуструктурированное интервью для оценки переживания «цифровой близости».
\end{itemize}

\subsection{Дизайн эксперимента}
\begin{itemize}
\item Контрольная группа (общение со стандартной, «безличной» версией бота).
\item Экспериментальная группа (общение с глубоко кастомизированным ботом — имя, аватар, стиль общения, бэкстори).
\item Серии коммуникативных сессий (нейтральная тема / личная / стрессовая).
\end{itemize}

\newpage
\section{Результаты исследования и их интерпретация}

\subsection{Анализ самораскрытия}
Сравнение глубины и интимности раскрываемой информации в группах. Влияние визуального образа (аватара) на готовность делиться личным.
% Текст раздела.

\subsection{Динамика регуляции аффекта}
Изменение эмоционального состояния до и после сессий. Как персона ИИ помогает снижать тревожность или, наоборот, может вызывать дискомфорт.
% Текст раздела.

\subsection{Феномен «цифровой близости»}
Качественный анализ интервью: испытывали ли пользователи чувство понимания, привязанности, интимности? Роль кастомизации в создании иллюзии «живого» общения.
% Текст раздела.

\subsection{Зависимость от виртуального партнёра}
Анализ признаков эмоциональной привязанности и рисков формирования зависимого паттерна взаимодействия с ИИ-компаньоном.
% Текст раздела.

\subsection{Сравнение с реальными межличностными отношениями}
Восприятие пользователями различий и сходств между взаимодействием с кастомизированным ИИ и отношениями с реальным партнёром.
% Текст раздела.

\subsection{Качественные данные: отзывы пользователей}
Обобщение субъективных оценок и переживаний респондентов по результатам интервью и открытых вопросов.
% Текст раздела.

\subsection{Экспертная оценка: мнение специалиста (психолог-психиатр)}
Интерпретация результатов с точки зрения клинической и консультативной практики; риски и рекомендации.
% Текст раздела.

\subsection{Обсуждение результатов}
Подтверждение или опровержение гипотез. Неожиданные находки.
% Текст раздела.

\newpage
\section*{Заключение}
\addcontentsline{toc}{section}{Заключение}
\begin{itemize}
\item Основные выводы исследования.
\item Ответ на вопрос, поставленный во введении.
\item Практические рекомендации для разработчиков подобных систем.
\item Ограничения исследования и перспективы дальнейшей разработки темы.
\end{itemize}

\end{document}